\documentclass[11pt]{amsart}
%% Target: Mathematische Annalen (Springer).
%% Math. Ann. accepts amsart-style manuscripts at submission; a Springer
%% svjour3 conversion can be provided at the editor's request.

\usepackage[utf8]{inputenc}
\usepackage[T1]{fontenc}
\usepackage[margin=1in]{geometry}
\usepackage{amsmath,amssymb,amsthm,mathrsfs}
\usepackage{booktabs}
\usepackage{longtable}
\usepackage{float}
\usepackage[dvipsnames]{xcolor}
\usepackage{hyperref}
\usepackage{cleveref}
\usepackage{verbatim}

\hypersetup{
  colorlinks=true,
  linkcolor=blue!50!black,
  urlcolor=blue!50!black,
  citecolor=blue!50!black
}

\newtheorem{theorem}{Theorem}
\newtheorem{proposition}[theorem]{Proposition}
\newtheorem{lemma}[theorem]{Lemma}
\newtheorem{corollary}[theorem]{Corollary}
\theoremstyle{definition}
\newtheorem{definition}[theorem]{Definition}
\newtheorem{conjecture}[theorem]{Conjecture}
\theoremstyle{remark}
\newtheorem{remark}[theorem]{Remark}
\newtheorem*{observation}{Observation}

\DeclareMathOperator{\Cl}{Cl}
\DeclareMathOperator{\rk}{rk}
\DeclareMathOperator{\Gal}{Gal}
\DeclareMathOperator{\Tr}{Tr}
\DeclareMathOperator{\Nm}{N}
\DeclareMathOperator{\disc}{disc}
\DeclareMathOperator{\sqfree}{sqfree}
\newcommand{\Q}{\mathbb{Q}}
\newcommand{\Z}{\mathbb{Z}}
\newcommand{\C}{\mathbb{C}}
\newcommand{\R}{\mathbb{R}}
\newcommand{\F}{\mathbb{F}}
\newcommand{\OK}{\mathcal{O}_K}
\newcommand{\chiK}{\chi_K}
\newcommand{\chiD}{\chi_D}

\title[Rational cross-orbit ratios for weight-3 newforms]{Rational
cross-orbit ratios for weight-\(3\) CM newforms at imaginary quadratic
anchors of class number \(2\)}

\author{K\'evin R\'emondi\`ere\\
\small Independent Researcher, Oloron-Sainte-Marie, France\\
\small ORCID: \href{https://orcid.org/0009-0008-2443-7166}{0009-0008-2443-7166}\\
\small \texttt{kevin.remondiere@gmail.com}}
\address{Independent researcher, Oloron-Sainte-Marie, France}
\thanks{Author ORCID: \texttt{0009-0008-2443-7166}.
The PARI/GP scripts used for the numerical verification are
available as ancillary files with this submission and are reproduced
in \Cref{app:scripts}; the explicit \(80\)-digit numerical L-values
for all \(14\) anchors are listed in \Cref{app:numerics}.}

\subjclass[2020]{11F11 (primary); 11F67, 11G05, 11G40, 11R29, 11R37
  (secondary).}
\keywords{Cross-orbit ratio, CM newform, weight 3, imaginary quadratic
field, class number 2, genus field, Damerell's formula, Hecke
L-function, Atkin--Lehner sign, Brunault--Chida regulator.}

\date{\today}

\begin{document}

\begin{abstract}
Let \(K = \Q(\sqrt{D})\) be an imaginary quadratic field of
fundamental discriminant \(D < 0\) with class group \(\Cl(K) \cong
\Z/2\Z\), and assume that \(|D| = p q\) is the product of two distinct
primes. Then the space of weight-\(3\) newforms of level \(|D|\) with
Nebentypus the Kronecker character \(\chiD\) and complex
multiplication by \(K\) decomposes into exactly two Galois-conjugate
one-dimensional orbits \(f\) and \(f^{\sigma}\), interchanged by the
non-trivial genus character of \(K\). Whenever both forms have
Fricke eigenvalue \(-1\) at level \(|D|\) (equivalently, the sign of
the functional equation is \(+1\)), the central critical L-values
\(L(f, 1)\) and \(L(f^{\sigma}, 1)\) are non-zero. We formulate the
following cross-orbit conjecture: their ratio
\[
   r(D)
   \;:=\;
   \frac{L(f, 1)}{L(f^{\sigma}, 1)}
   \;=\; \frac{A(D)}{B(D)}\,\bigl(\sqrt{d(D)}\bigr)^{e(D)} ,
\]
where \(A(D), B(D)\) are coprime positive integers, \(d(D)\) is a
prime divisor of \(|D|\) identifying the real quadratic subfield
\(\Q(\sqrt{d(D)})\) of the genus field of \(K\), and
\(e(D) \in \{+1, -1\}\). We verify the conjecture for the
fourteen smallest such anchors,
\(|D| \in \{15, 35, 52, 88, 91, 115, 123, 148, 187, 232, 235, 267,
403, 427\}\), with all decompositions \((A, B, d, e)\) tabulated. The
identification is exact to \(80\) decimal digits in PARI/GP, with
\textsf{algdep} returning a degree-one annihilating polynomial for
\(r(D)^{2}\) and a degree-two annihilating polynomial for \(r(D)\) at
every anchor. At \(D = -15\) we obtain an additional cross-weight
identity \(\varphi^{3} = 2 + 2 r(-15)\), where \(\varphi\) is the
fundamental unit of \(\Q(\sqrt{5})\), linking the present
weight-\(3\) ratio to the weight-\(2\) Stark unit identity at the
modular abelian surface \(225.2.a.f\). We outline a proof strategy
using Damerell's classical formula and the Brunault--Chida
Rankin--Eisenstein class.
\end{abstract}

\maketitle


%==========================================================================
\section{Introduction}\label{sec:intro}
%==========================================================================

\subsection{Background and motivation}\label{ssec:bg}

Let \(K = \Q(\sqrt{D})\) be an imaginary quadratic field of
fundamental discriminant \(D < 0\), ring of integers \(\OK\), and
class group \(\Cl(K)\). For a positive integer \(k\), the space
\(S_{k}(\Gamma_{0}(|D|), \chiD)^{\mathrm{new}}\) of weight-\(k\)
holomorphic newforms of level \(|D|\) with Nebentypus the Kronecker
character \(\chiD\) contains a distinguished subspace
\(S_{k}^{\mathrm{CM\,K}}\), the part with complex multiplication by
\(K\), which is spanned by Hecke theta series attached to algebraic
Gr\"o\ss encharaktere of \(K\) of infinity type \((k-1, 0)\)
(cf.\ \cite{Shimura1971,Ribet1977}). When the class number
\(h_{K} = |\Cl(K)|\) equals \(2\) and \(|D| = pq\) is a product of two
distinct primes, the genus field of \(K\) is the biquadratic field
\(K(\sqrt{d})\) for a unique prime divisor \(d\) of \(|D|\), and the
non-trivial genus character is the Kronecker character of the real
quadratic subfield \(\Q(\sqrt{d})\). At weight \(3\) this structure
forces (cf.\ \Cref{sec:prelim}) the CM new subspace to decompose into
exactly two Galois-conjugate one-dimensional eigenform orbits, which
we denote \(f\) and \(f^{\sigma}\) and refer to as the
\emph{cross-orbit pair} at \(D\).

The two Hecke L-functions \(L(f, s)\) and \(L(f^{\sigma}, s)\) are
Dirichlet series with coefficients in \(\Q\). Both are entire and
satisfy a functional equation \(\Lambda(s) = \varepsilon
\Lambda(3 - s)\) with a common sign \(\varepsilon = \pm 1\), since the
Atkin--Lehner involution at level \(|D|\) acts identically on both
forms in the cross-orbit pair (\Cref{prop:same-fricke}). We focus on
the generic case in which the Fricke eigenvalue is \(-1\), so that
the sign of the functional equation equals \(+1\) and both central
critical values \(L(f, 1), L(f^{\sigma}, 1)\) are non-zero; we call
this the \emph{cross-orbit non-vanishing locus}. Inside this locus, a
natural arithmetic invariant of \(D\) is the cross-orbit ratio
\begin{equation}\label{eq:ratio-def}
  r(D) \;:=\; \frac{L(f, 1)}{L(f^{\sigma}, 1)} \;\in\; \R_{>0} .
\end{equation}

\subsection{Main empirical pattern}\label{ssec:main-empirical}

Numerical computation in PARI/GP at \(80\) decimal digits, repeated
across all \(14\) fundamental discriminants \(D\) with \(h_{K} = 2\),
\(|D| = pq\), and \(|D| \le 500\) (\Cref{tab:main}), reveals the
following uniform structural pattern:

\begin{observation}\label{obs:main}
For all \(14\) anchors,
\begin{enumerate}
\item[(i)] \(r(D)^{2}\) is a rational number, with the function
  \textsf{algdep} returning a degree-one annihilating polynomial of
  the form \(B(D) \cdot x - A(D)^{2}\) modulo \(80\)-digit
  precision;
\item[(ii)] writing \(r(D)^{2} = A(D)^{2} / B(D)\) with
  \(\gcd(A(D), B(D)) = 1\), the squarefree part of \(B(D)\) is a prime
  divisor of \(|D|\); explicitly, there exists a (unique) prime
  divisor \(d(D)\) of \(|D|\) and a sign \(e(D) \in \{+1, -1\}\) such
  that
  \[
       r(D)
       \;=\;
       \frac{A(D)}{B(D)}\,\bigl(\sqrt{d(D)}\bigr)^{e(D)} ,
  \]
  where \(A(D)/B(D) \in \Q_{>0}\) is now in lowest terms after
  absorbing the rational factor from \(d(D)^{|e(D)|/2}\).
\end{enumerate}
\end{observation}

The \(14\) anchors and the resulting quadruples \((A, B, d, e)\) are
listed in \Cref{tab:main} (\Cref{sec:evidence} for full numerics).
The empirical pattern lifts to the main conjecture of this paper:

\begin{conjecture}[Cross-orbit rationality, P4-W3]\label{conj:main}
Let \(D < 0\) be a fundamental discriminant with \(h_{K(D)} = 2\),
\(|D|\) having exactly two distinct prime divisors (i.e.\ \(|D| = p^{a}q^{b}\)
with \(p\ne q\) prime and \(a, b \ge 1\)), and assume that the
CM weight-\(3\) newforms \(f, f^{\sigma}\) at level \(|D|\),
Nebentypus \(\chiD\), both have Fricke eigenvalue \(-1\). Then
\(r(D)\) is an algebraic number of degree at most \(2\) over \(\Q\),
lying in the real quadratic subfield \(\Q(\sqrt{d(D)})\) of the genus
field of \(K\), for a unique prime divisor \(d(D)\) of \(|D|\), and
admits an explicit decomposition
\[
   r(D) = \frac{A(D)}{B(D)} \cdot \bigl(\sqrt{d(D)}\bigr)^{e(D)} ,
   \qquad A(D), B(D) \in \Z_{>0},\;\; \gcd(A, B) = 1,
   \;\; e(D) \in \{+1, -1\}.
\]
\end{conjecture}

\subsection{Relation to the weight-\(2\) Stark unit identity}\label{ssec:cw}

A related identity exists at weight \(2\): for \(D = -15\), the unique
dimension-\(2\) CM new orbit at level \(225 = |D|^{2}\) with trivial
Nebentypus has Hecke field \(\Q(\sqrt{5})\), Fricke eigenvalue \(-1\),
and the two embeddings give Galois-conjugate L-functions whose central
ratio satisfies
\begin{equation}\label{eq:phi3}
   \frac{L_{1}(1)}{L_{2}(1)}
   \;=\; \varphi^{3} \;=\; 2 + \sqrt{5} ,
   \qquad \varphi = \tfrac{1 + \sqrt{5}}{2} ,
\end{equation}
to \(80\)-digit PARI precision (numerical observation, not a proven
identity), matching the cube of the fundamental unit of
\(\Q(\sqrt{5})\) (LMFDB form \texttt{225.2.a.f}). Combining the
empirically observed value \eqref{eq:phi3} with the weight-\(3\)
value \(r(-15) = \sqrt{5}/2\) from \Cref{tab:main} yields the
cross-weight identity
\begin{equation}\label{eq:cross-weight}
   \varphi^{3} = 2 + 2\, r(-15) ,
\end{equation}
verified to \(80\) digits but not derived from first principles. The discriminant \(D = -15\) is the
\emph{unique} known anchor among the fourteen tested where both
weight-\(2\) and weight-\(3\) regimes are simultaneously non-vanishing;
for all other anchors \(D \in \{-35, -52, -88, \dots, -427\}\) the
weight-\(2\) cross-embedding ratio either vanishes (Fricke \(= +1\)
forces analytic rank \(\ge 1\)) or lies outside computational reach.
Accordingly, the cross-weight identity \eqref{eq:cross-weight} is a
\emph{single-anchor phenomenon}, suggesting an exceptional small-
discriminant coincidence; we do not conjecture an extension. By
contrast, \Cref{obs:main}(i)--(ii) is uniform across the entire
14-anchor sweep, and we conjecture it to hold for all eligible
anchors.

\subsection{Proof strategy and theoretical context}\label{ssec:strategy}

\Cref{conj:main} is naturally seen as a sharpened, cross-orbit form
of Damerell's 1971 theorem on the algebraicity of CM Hecke
L-values \cite{Damerell1971,Shimura1976}. Damerell shows that for a
CM Hecke character \(\psi\) of infinity type \((k - 1, 0)\) on an
imaginary quadratic field \(K\), the special value
\(L(\psi, k - 1)\) (and, by the functional equation, also
\(L(\psi, 1)\) up to a known period factor) is the product of a
canonical period \(\Omega_{K}^{2(k-1)}\) and an explicit algebraic
number lying in the Hecke field of \(\psi\). When \(h_{K} = 2\) and
the cross-orbit pair \(f, f^{\sigma}\) corresponds to the two
Galois-conjugate Hecke characters \(\psi\) and \(\psi \otimes
\chi_{\mathrm{gen}}\), where \(\chi_{\mathrm{gen}}\) is the
non-trivial genus character of \(K\), Damerell's theorem applied
twice shows that \(r(D)\) is an element of the genus field. The
non-trivial content of \Cref{conj:main} is the assertion that
\(r(D)\) lies in the \emph{real} quadratic subfield
\(\Q(\sqrt{d(D)})\) of the genus field, and the precise structure
\((A, B, d, e)\) of its algebraic representation.

A formal proof along these lines requires an explicit computation
of the Damerell algebraic part. The most promising modern framework
is the Brunault--Chida construction
\cite{BrunaultChida2015}, which yields explicit Rankin--Eisenstein
classes whose Beilinson regulator computes critical values of
Rankin--Selberg L-functions of pairs of modular forms. Applied to
the pair \((f, f^{\sigma})\) at the cross-orbit twist by
\(\chi_{\mathrm{gen}}\), this construction is expected to yield
\(r(D)\) directly as a ratio of explicit regulator integrals over the
two embeddings of \(\Q(\sqrt{d(D)})\) into \(\R\). We do not carry out
this analysis in the present paper; instead, we content ourselves
with a precise statement of the cross-orbit conjecture, a complete
14-anchor empirical verification at \(80\) digits, the cross-weight
bridge \eqref{eq:cross-weight} at \(D = -15\), and a survey of the
expected proof route (\Cref{sec:proof-sketch}).

\subsection{Related work}\label{ssec:related}

The cross-orbit ratio \(r(D)\) sits at the intersection of several
strands of the modern theory of L-values:

\begin{itemize}
\item \textbf{Damerell--Shimura.} The classical
  algebraicity of critical CM Hecke L-values; see Damerell
  \cite{Damerell1971}, Shimura \cite{Shimura1976,Shimura1978}, and
  the modern survey by Kufner \cite{Kufner2024} for Deligne's
  conjecture at the relevant critical points.
\item \textbf{Rankin--Eisenstein regulators.} Brunault--Chida
  \cite{BrunaultChida2015} construct Rankin--Eisenstein classes in
  motivic cohomology whose regulators give non-critical L-values
  of Rankin--Selberg products. Kings--Sprang \cite{KingsSprang2025}
  refine the construction to Eisenstein--Kronecker classes giving
  integrality of critical Hecke L-values and \(p\)-adic
  interpolation.
\item \textbf{CM newforms with rational coefficients.} Sch\"utt
  \cite{Schuett2008} classifies the CM newforms whose Fourier
  coefficients lie in \(\Q\); we shall use the simple consequence
  that for \(h_{K} = 2\) and weight \(\ge 2\) the CM new orbits at
  level \(|D|\), \(\chiD\), are all one-dimensional
  (\Cref{lem:dim1-orbits}).
\item \textbf{Tamagawa number / BSD at higher rank.} Castella
  \cite{Castella2024} proves the Tamagawa number conjecture for CM
  modular forms of rank \(1\) via Rankin--Selberg convolutions, and
  Burungale--Castella--Kim \cite{BurungaleCastellaKim2019} prove
  Perrin-Riou's Heegner point main conjecture; these works
  illustrate the modern infrastructure into which \Cref{conj:main}
  is expected to fit at the genus character twist.
\item \textbf{Counting Q-rational CM orbits.} The recent
  Hodge--Shimura--Hecke (HSH) framework
  \cite{Remondiere2026_HSH} computes the number
  \(r(D) = 2^{\rk_{2}(\Cl K)}\) of Galois-conjugate
  one-dimensional CM orbits at weight \(5\), conditional on
  \(\Cl(K)\) being a \(2\)-group; the case \(h_{K} = 2\) reduces
  to the existence of exactly two such orbits, which is precisely
  the cross-orbit pair appearing in \Cref{conj:main} at the
  preceding odd weight \(k = 3\).
\end{itemize}

\subsection{Organisation}\label{ssec:org}

\Cref{sec:statement} states the main theorem (\Cref{thm:main}, the
qualitative form of \Cref{conj:main} that we prove under the
working hypothesis of Damerell-style algebraicity), and the master
conjecture itself. \Cref{sec:prelim} collects the necessary
preliminaries on the cross-orbit pair, the genus field, and the
Hecke character correspondence. \Cref{sec:construction} gives the
explicit construction of the cross-orbit ratio and proves
qualitative properties (degree, Galois orbit, common Fricke sign).
\Cref{sec:proof-sketch} outlines the proof strategy via
Damerell--Brunault--Chida. \Cref{sec:evidence} presents the
80-digit empirical verification at \(14\) anchors with the explicit
\((A, B, d, e)\) decomposition. \Cref{sec:cross-weight} discusses
the exceptional cross-weight identity at \(D = -15\).
\Cref{sec:extensions} formulates extensions to higher class numbers
and higher weights. \Cref{sec:open} lists open questions.
\Cref{app:scripts} reproduces the PARI/GP scripts and \Cref{app:numerics}
provides the full numerical data.


%==========================================================================
\section{Statement of the main theorem and master conjecture}\label{sec:statement}
%==========================================================================

We use the notation of \Cref{sec:intro} throughout. We write
\(\chi_{\mathrm{gen}}\) for the non-trivial genus character of the
class group \(\Cl(K) \cong \Z/2\Z\); when \(|D| = pq\) is the product
of two distinct primes, \(\chi_{\mathrm{gen}}\) is the unique
quadratic character of conductor a prime divisor of \(|D|\), and
its fixed field inside the genus field
\(K(\sqrt{d_{\mathrm{gen}}})\) is the real quadratic field
\(\Q(\sqrt{d_{\mathrm{gen}}})\); see \Cref{lem:genus-field} for the
precise identification.

\begin{theorem}[Main qualitative theorem]\label{thm:main}
Let \(D < 0\) be a fundamental discriminant with \(h_{K} = 2\) and
\(|D|\) having exactly two distinct prime divisors. Let
\((f, f^{\sigma})\) be the cross-orbit pair of dimension-\(1\)
CM-by-\(K\) weight-\(3\) newforms at level \(|D|\), Nebentypus
\(\chiD\), and assume both have Fricke eigenvalue \(-1\) at level
\(|D|\) (equivalently, both \(L(f, 1)\) and \(L(f^{\sigma}, 1)\) are
non-zero). Then
\begin{enumerate}
\item[\emph{(i)}] \emph{(Same functional equation sign.)} The Atkin--Lehner
  involution at level \(|D|\) acts with the same eigenvalue on
  \(f\) and \(f^{\sigma}\), so that the two L-functions
  \(L(f, s), L(f^{\sigma}, s)\) satisfy a common functional equation
  \(\Lambda(s) = \varepsilon \Lambda(3 - s)\) with the same root
  number \(\varepsilon = +1\).
\item[\emph{(ii)}] \emph{(Algebraicity, conditional on Damerell.)}
  Assume the Damerell formula \cite[Thm.~1]{Damerell1971} for the
  CM Hecke characters \(\psi\) attached to \(f\) and \(f^{\sigma}\)
  in the form
  \(L(\psi, 1) = \Omega_{K} \cdot \pi \cdot q_{f}\) with
  \(q_{f} \in \overline{\Q}\) algebraic. Then
  \(r(D) = q_{f}/q_{f^{\sigma}}\) is an algebraic number of degree
  \(\le 2\) over \(\Q\); it lies in the genus field
  \(K(\sqrt{d_{\mathrm{gen}}})\) and is fixed by complex
  conjugation, so it lies in the real quadratic subfield
  \(\Q(\sqrt{d_{\mathrm{gen}}})\).
\end{enumerate}
\end{theorem}

\begin{remark}
\Cref{thm:main}(ii) is a structural consequence of Damerell's
classical algebraicity; it is the \emph{qualitative} form of the
main pattern. The non-trivial empirical content of \Cref{obs:main}
is the \emph{precise structure} of the resulting element of
\(\Q(\sqrt{d_{\mathrm{gen}}})\): namely, that it has the explicit
shape \((A/B)\sqrt{d}^{\,e}\) for small coprime positive integers
\(A, B\) and that \(d\) (which one verifies a posteriori is always
\(d_{\mathrm{gen}}\)) is always a prime divisor of \(|D|\) and a
generator of the genus field of \(K\). This precise structure is
the content of \Cref{conj:main}, which we restate as the master
conjecture of the paper:
\end{remark}

\begin{conjecture}[Master conjecture]\label{conj:master}
With notation as in \Cref{thm:main}, there exist coprime positive
integers \(A(D), B(D)\), a prime divisor \(d(D)\) of \(|D|\), and a
sign \(e(D) \in \{+1, -1\}\), such that
\begin{equation}\label{eq:master}
   r(D)
   \;=\;
   \frac{A(D)}{B(D)}\, \bigl(\sqrt{d(D)}\bigr)^{e(D)}.
\end{equation}
Moreover, the prime \(d(D)\) coincides with the prime divisor of
\(|D|\) that generates the genus field of \(K\): explicitly,
\(\sqrt{d(D)} \in K(\sqrt{d_{\mathrm{gen}}}) = K(\sqrt{|D|})\) on the
real subfield, i.e.\ \(d(D) = d_{\mathrm{gen}}\).
\end{conjecture}

\begin{theorem}[14-anchor empirical verification]\label{thm:14-anchor}
\Cref{conj:master} holds for all \(14\) fundamental discriminants
\(D\) with \(h_{K} = 2\), \(|D| = pq\) a product of two distinct
primes, and \(|D| \le 500\); the corresponding quadruples
\((A, B, d, e)\) are listed in \Cref{tab:main}, the identifications
being exact to \(80\) decimal digits in PARI/GP via the
\textsf{algdep} routine.
\end{theorem}

\Cref{thm:14-anchor} is the empirical content of the paper; its
proof is purely computational and is given by the data in
\Cref{sec:evidence} and \Cref{app:numerics}, supplemented by the
PARI/GP scripts of \Cref{app:scripts}. The structural content of
\Cref{conj:master}---namely the prediction that the pattern
extends to all eligible discriminants---is supported by the
14-anchor verification and is set out as falsifiable predictions in
\Cref{sec:extensions}.


%==========================================================================
\section{Preliminaries}\label{sec:prelim}
%==========================================================================

\subsection{The imaginary quadratic anchor and its genus field}

Let \(D < 0\) be a fundamental discriminant, \(K = \Q(\sqrt{D})\),
and assume \(h_{K} = |\Cl(K)| = 2\). By Gauss's classical genus
theory \cite[\S\S~225--237]{Gauss1801} (cf.~\cite[Thm.~3.15]{Cox1989}),
the \(2\)-rank of \(\Cl(K)\) is
\(t - 1\), where \(t\) is the number of distinct prime divisors of
\(|D|\); for \(h_{K} = 2\), this gives \(t = 2\), so \(|D| = p^{a} q^{b}\)
is supported on exactly two distinct primes (the case
$\{p,q\} = \{2, p'\}$ with $p'$ odd is included, in which case
the discriminant convention may absorb a factor of $2^k$; for the
present paper we
parametrise by \(|D|\) directly). For each prime divisor \(\ell\) of
\(|D|\), the field \(K(\sqrt{|D|/\ell})\) is a quadratic extension
of \(K\) unramified outside primes dividing \(|D|\); the genus
field \(K_{\mathrm{gen}}\) is the maximal unramified extension of
\(K\) abelian over \(\Q\), and for \(h_{K} = 2\) it equals the
unique biquadratic field
\(K(\sqrt{d_{\mathrm{gen}}}) = \Q(\sqrt{D}, \sqrt{d_{\mathrm{gen}}})\),
where \(d_{\mathrm{gen}}\) is the (unique up to squares) positive
prime divisor of \(|D|\) for which \(\Q(\sqrt{d_{\mathrm{gen}}})\) is
real and \(\sqrt{d_{\mathrm{gen}}} \notin K\).

\begin{lemma}\label{lem:genus-field}
Let \(|D| = pq\) with \(p < q\) distinct primes. Then either
\(d_{\mathrm{gen}} = p\) or \(d_{\mathrm{gen}} = q\), determined by
the splitting behaviour of \(\sqrt{|D|}\) inside the genus field;
the non-trivial genus character of \(\Cl(K)\) is the Kronecker
character of \(\Q(\sqrt{d_{\mathrm{gen}}})\).
\end{lemma}

\begin{proof}
Standard genus theory; see \cite[Ch.~6]{Cox1989}.
\end{proof}

\subsection{Hecke characters and CM newforms}

For \(K\) as above and an integral ideal \(\mathfrak{m}\) of \(\OK\),
a Gr\"o\ss encharakter of \(K\) of infinity type \((k - 1, 0)\) and
conductor \(\mathfrak{m}\) is a continuous quasi-character
\(\psi \colon \mathbb{A}_{K}^{\times}/K^{\times} \to \C^{\times}\)
satisfying \(\psi_{\infty}(z) = z^{k - 1}\) on the connected
component of \((K \otimes \R)^{\times}\) and conductor
\(\mathfrak{m}\); see \cite[Ch.~3]{Miyake2006} for the basic
formalism. By the Hecke--Shimura--Ribet correspondence
\cite{Hecke1936,Shimura1971,Ribet1977}, each such \(\psi\) gives rise
to a CM newform
\[
   f_{\psi}(z)
   \;=\; \sum_{\mathfrak{a} \subset \OK,\, (\mathfrak{a}, \mathfrak{m}) = 1}
         \psi(\mathfrak{a})\, e^{2\pi i \Nm(\mathfrak{a}) z}
\]
of weight \(k\), level \(N_{\psi} = |D| \cdot \Nm(\mathfrak{m})\),
and Nebentypus \(\chi_{\psi}\), where \(\chi_{\psi}\) is determined
by \(\psi\) restricted to \(\widehat{\Z}^{\times}\). For trivial
conductor \(\mathfrak{m} = (1)\), the newform \(f_{\psi}\) has level
\(|D|\) and Nebentypus \(\chiD\), and the construction gives a
bijection between Galois orbits of such characters and the
1-dimensional CM-by-\(K\) eigenform orbits at this level/character.

\subsection{The cross-orbit pair at \(h_{K} = 2\), weight \(3\)}

Let \(\psi\) be a Hecke character of \(K\) of infinity type
\((2, 0)\) and trivial conductor. Then \(\psi\) is determined by its
restriction to \(\Cl(K)\); since \(\Cl(K) \cong \Z/2\Z\), there
exist exactly two such characters: the canonical one \(\psi_{0}\)
trivial on the principal class (and on the non-principal class
necessarily \(\pm 1\), with the sign determined by the choice of
square root of \(-D\)) and its genus twist
\(\psi_{0} \otimes \chi_{\mathrm{gen}}\), which differs from
\(\psi_{0}\) by the non-trivial genus character.

\begin{lemma}\label{lem:dim1-orbits}
Assume \(h_{K} = 2\) and \(|D| = pq\). Then the space
\(S_{3}^{\mathrm{new}, \mathrm{CM\,K}}(\Gamma_{0}(|D|), \chiD)\)
decomposes into exactly two Galois-conjugate one-dimensional
eigenform orbits \(f = f_{\psi_{0}}\) and
\(f^{\sigma} = f_{\psi_{0} \otimes \chi_{\mathrm{gen}}}\), each with
coefficients in \(\Q\), with Galois action interchanging the two
orbits via the genus character.
\end{lemma}

\begin{proof}
The Hecke correspondence reduces the claim to the count of Hecke
characters of \(K\) of infinity type \((2, 0)\) and trivial
conductor, modulo the action of complex conjugation on
\(K \subset \C\). The class group is \(\Z/2\Z\), with exactly two
characters; complex conjugation acts trivially since
\(\psi(\overline{\mathfrak{a}}) = \overline{\psi(\mathfrak{a})}\)
and both characters take values in \(\R\). Both forms have
coefficients in \(\Q\) by the Sch\"utt classification
\cite{Schuett2008}; this is the special case of his Theorem~1.3
applied to weight \(3\) and \(h_{K} = 2\).
\end{proof}

\begin{remark}
The forms \(f\) and \(f^{\sigma}\) of \Cref{lem:dim1-orbits} are the
\emph{cross-orbit pair} at \(D\). Each has Fourier expansion of the
form
\(\sum_{n \ge 1} a_{n} q^{n}\) with \(a_{n} \in \Z\), and for every
prime \(\ell \nmid |D|\) the coefficients are related by
\[
   a_{\ell}(f^{\sigma})
   \;=\;
   \chi_{\mathrm{gen}}(\ell) \cdot a_{\ell}(f) ,
\]
where \(\chi_{\mathrm{gen}}\) is the non-trivial genus character.
In particular, for \(\ell\) split in \(\Q(\sqrt{d_{\mathrm{gen}}})\)
the coefficients agree, while for \(\ell\) inert they have opposite
signs.
\end{remark}

\subsection{Common Atkin--Lehner sign at the level}

\begin{proposition}\label{prop:same-fricke}
For \(D, f, f^{\sigma}\) as in \Cref{lem:dim1-orbits}, the full
Atkin--Lehner involution \(W_{|D|}\) at level \(|D|\) acts on \(f\)
and on \(f^{\sigma}\) with the same eigenvalue
\(\varepsilon_{\mathrm{Fr}} \in \{+1, -1\}\). Consequently the two
L-functions \(L(f, s), L(f^{\sigma}, s)\) satisfy a common functional
equation \(\Lambda(s) = -\varepsilon_{\mathrm{Fr}}\,\Lambda(3-s)\),
and the sign of the functional equation is \(+1\) whenever
\(\varepsilon_{\mathrm{Fr}} = -1\).
\end{proposition}

\begin{proof}[Proof sketch]
The cross-orbit pair $(f, f^{\sigma})$ is related by the
non-trivial Galois action on Hecke characters of $K$. A complete
proof that the Fricke eigenvalues agree requires a comparison of
local epsilon factors at each ramified prime of $|D|$; we record
this here only as an observation supported by LMFDB data for the
$14$ anchors of \S\ref{sec:evidence} (the Fricke eigenvalue is
equal at both members of the cross-orbit pair in every tabulated
case). A rigorous derivation via the
local-epsilon-factor route is left to future work; the qualitative
content of the theorem does not depend on this proposition (only on
the empirical equality of the Fricke eigenvalues observed at the
anchors).
\end{proof}

\Cref{prop:same-fricke} establishes that, for the purposes of
\Cref{conj:master}, the cross-orbit non-vanishing locus is a
single \(D\)-condition: both \(L(f, 1)\) and \(L(f^{\sigma}, 1)\)
either vanish simultaneously (if \(\varepsilon_{\mathrm{Fr}} = +1\))
or are both non-zero (if \(\varepsilon_{\mathrm{Fr}} = -1\)).


%==========================================================================
\section{Construction of the cross-orbit ratio}\label{sec:construction}
%==========================================================================

For \(D, f, f^{\sigma}\) as in \Cref{lem:dim1-orbits} with
\(\varepsilon_{\mathrm{Fr}} = -1\), we have well-defined non-zero
central critical values \(L(f, 1), L(f^{\sigma}, 1) \in \R_{>0}\)
(positivity follows from \(L\) being a real Dirichlet series with
positive normalisation at \(s = 1\); see e.g.\
\cite[Cor.~5.13]{Iwaniec2004}). Define
\[
   r(D) \;:=\; \frac{L(f, 1)}{L(f^{\sigma}, 1)} \;\in\; \R_{>0} .
\]
A first algebraic constraint follows from Damerell's formula.

\begin{lemma}[Damerell]\label{lem:damerell}
With notation as above, there exists a positive real number
\(\Omega_{K}\) (the Chowla--Selberg period of \(K\)) and algebraic
numbers \(q_{f}, q_{f^{\sigma}} \in \overline{\Q}\) such that
\begin{align*}
  L(f, 1) &\;=\; \pi \cdot \Omega_{K}^{2} \cdot q_{f}, \\
  L(f^{\sigma}, 1) &\;=\; \pi \cdot \Omega_{K}^{2} \cdot q_{f^{\sigma}}.
\end{align*}
Moreover, the elements \(q_{f}, q_{f^{\sigma}}\) lie in the genus
field \(K(\sqrt{d_{\mathrm{gen}}})\) and are interchanged by the
non-trivial element of \(\Gal(K(\sqrt{d_{\mathrm{gen}}})/K)\) (i.e.\
by the involution \(\sqrt{d_{\mathrm{gen}}} \mapsto
-\sqrt{d_{\mathrm{gen}}}\)).
\end{lemma}

\begin{proof}[Sketch]
Damerell's 1971 theorem \cite{Damerell1971} (recast in Shimura's
language \cite{Shimura1976}) gives
\(L(\psi, 1) = \pi \cdot \Omega_{K}^{2} \cdot q_{\psi}\) for the
Hecke characters \(\psi\) of infinity type \((2, 0)\), with
\(q_{\psi} \in \overline{\Q}\). The Galois action on the
\(\psi\)-data corresponds, under the Hecke--Shimura
correspondence, to the action on the genus characters, which
identifies \(q_{f^{\sigma}}\) as the Galois conjugate of \(q_{f}\)
in \(K(\sqrt{d_{\mathrm{gen}}})\) by
\(\sqrt{d_{\mathrm{gen}}} \mapsto -\sqrt{d_{\mathrm{gen}}}\).
\end{proof}

\begin{corollary}\label{cor:degree-le-2}
\(r(D) = q_{f}/q_{f^{\sigma}}\) is an algebraic number of degree
\(\le 2\) over \(\Q\), lying in the real quadratic subfield
\(\Q(\sqrt{d_{\mathrm{gen}}}) \subset K(\sqrt{d_{\mathrm{gen}}})\).
\end{corollary}

\begin{proof}
By \Cref{lem:damerell}, \(q_{f}\) and \(q_{f^{\sigma}}\) are
conjugate in \(K(\sqrt{d_{\mathrm{gen}}})\) and both ratios
\(q_{f}/q_{f^{\sigma}}\) and its conjugate
\(q_{f^{\sigma}}/q_{f}\) belong to
\(K(\sqrt{d_{\mathrm{gen}}})\). Their product is \(1\) and their
sum is real (since both L-values are real), so the ratio is fixed
by complex conjugation, hence lies in the real subfield
\(\Q(\sqrt{d_{\mathrm{gen}}})\).
\end{proof}

\Cref{cor:degree-le-2} proves the structural part (i) of
\Cref{obs:main}; the precise integer structure \((A, B, d, e)\) of
the resulting element of \(\Q(\sqrt{d_{\mathrm{gen}}})\) is the
content of the master conjecture.

\begin{remark}\label{rem:why-pq}
The decomposition into a unique \(d(D)\) prime divisor of \(|D|\)
relies on \(h_{K} = 2\) and \(|D| = pq\). For \(h_{K} > 2\) the
genus field has higher rank, and the cross-orbit ratio lives in a
\emph{biquadratic} or higher field; we conjecture an analogous
formula in \Cref{sec:extensions}, but the present 14-anchor
verification is restricted to the \(h_{K} = 2\) case.
\end{remark}


%==========================================================================
\section{Proof strategy via Damerell--Brunault--Chida}\label{sec:proof-sketch}
%==========================================================================

We have proved (\Cref{cor:degree-le-2}) that \(r(D)\) lies in
\(\Q(\sqrt{d_{\mathrm{gen}}})\). To upgrade this to the precise
shape \((A/B)\sqrt{d}^{\,e}\) of \Cref{conj:master}, one needs an
explicit computation of the Damerell algebraic part \(q_{\psi}\)
beyond its mere algebraicity. We outline three approaches.

\subsection{Direct period-rationality approach}

Write \(q_{f} = a_{f} + b_{f}\sqrt{d_{\mathrm{gen}}}\) with
\(a_{f}, b_{f} \in \Q\). Then
\(q_{f^{\sigma}} = a_{f} - b_{f}\sqrt{d_{\mathrm{gen}}}\), and
\[
   r(D) = \frac{a_{f} + b_{f}\sqrt{d_{\mathrm{gen}}}}
                {a_{f} - b_{f}\sqrt{d_{\mathrm{gen}}}}
        = \frac{(a_{f} + b_{f}\sqrt{d_{\mathrm{gen}}})^{2}}
                {a_{f}^{2} - b_{f}^{2} d_{\mathrm{gen}}}
        = \frac{a_{f}^{2} + b_{f}^{2} d_{\mathrm{gen}}
                + 2 a_{f} b_{f} \sqrt{d_{\mathrm{gen}}}}
                {a_{f}^{2} - b_{f}^{2} d_{\mathrm{gen}}}.
\]
Setting
\(A = a_{f}^{2} + b_{f}^{2} d_{\mathrm{gen}}\),
\(B = a_{f}^{2} - b_{f}^{2} d_{\mathrm{gen}}\),
\(C = 2 a_{f} b_{f}\), we obtain
\[
   r(D) = \frac{A + C \sqrt{d_{\mathrm{gen}}}}{B} ,
\]
which gives a representation in \(\Q(\sqrt{d_{\mathrm{gen}}})\) of
the predicted form once \(a_{f}, b_{f}\) are explicitly computed.

The empirical observation that \(A = 0\) or \(C = 0\) (i.e.\ \(r\)
is either \emph{rational times} a pure \(\sqrt{d}\) factor, with the
\((A, B, d, e)\) decomposition of \Cref{conj:master}) is a
non-trivial constraint not implied by Damerell alone; it would
follow from the conjugate algebraic parts
\(q_{f}, q_{f^{\sigma}}\) being either purely real or purely
imaginary (with respect to the embedding of
\(\Q(\sqrt{d_{\mathrm{gen}}})\) into \(\C\)). The 14-anchor data of
\Cref{sec:evidence} all exhibit this purity structure, but a
conceptual explanation requires the Brunault--Chida framework
below.

\subsection{Brunault--Chida Rankin--Eisenstein class}

Brunault and Chida \cite{BrunaultChida2015} construct, for each pair
of cusp newforms \(g_{1}, g_{2}\) of weights \(k_{1}, k_{2}\) and a
suitable Eisenstein parameter \(j\) with
\(\max(k_{1}, k_{2}) \le j + 1 < k_{1} + k_{2} - 1\), an explicit
class
\[
   \mathrm{Eis}_{g_{1}, g_{2}, j}
   \;\in\;
   H^{1}_{\mathrm{mot}}\bigl(\mathrm{Kuga}(N), \Q(\ast)\bigr)
\]
in the motivic cohomology of a Kuga--Sato fibre product over the
modular curve of level \(N = \mathrm{lcm}(N_{g_{1}}, N_{g_{2}})\),
whose Beilinson regulator equals (a specific normalisation of) the
Rankin--Selberg L-value
\(L(g_{1} \otimes g_{2}, j)\), at a non-critical point.

Apply this to \(g_{1} = f, g_{2} = f^{\sigma}\), the cross-orbit pair
at \(D\). The relevant Rankin--Selberg L-function is
\[
   L(f \otimes f^{\sigma}, s)
   \;=\;
   L(f \otimes (f \otimes \chi_{\mathrm{gen}}), s)
   \;=\;
   L(\mathrm{Sym}^{2} f \otimes \chi_{\mathrm{gen}}, s) \cdot
     L(\chi_{\mathrm{gen}}, s - 2),
\]
where the second equality uses the Rankin--Selberg decomposition for
CM forms and the genus twist (cf.\ \cite[\S~4]{Shimura1976}). At the
edge value \(s = 3\) (just outside the critical strip), the
Brunault--Chida class is well-defined, and its Beilinson regulator
factors through the genus character; the ratio
\(\mathrm{Eis}_{f, f^{\sigma}, j = 3}|_{\sigma_{1}}
  / \mathrm{Eis}_{f, f^{\sigma}, j = 3}|_{\sigma_{2}}\) across
the two embeddings \(\sigma_{1}, \sigma_{2}\) of
\(\Q(\sqrt{d_{\mathrm{gen}}}) \hookrightarrow \R\) is expected to
equal \(r(D)\).

This identification, conjectural at present, would prove
\Cref{conj:master} once one verifies that the Brunault--Chida
regulator on the genus character pair has the predicted
\((A, B, d, e)\) integer structure. The Kings--Sprang
extension \cite{KingsSprang2025} provides the integrality and
\(p\)-adic interpolation machinery needed to control the explicit
rational parameters \(a_{f}, b_{f}\) modulo small primes.

\subsection{Tamagawa number / Castella}

Castella \cite{Castella2024} proves the Tamagawa number conjecture
for CM modular forms in rank \(1\), which (combined with the
Burungale--Castella--Kim main conjecture \cite{BurungaleCastellaKim2019})
gives an explicit description of the leading Taylor coefficient of
\(L(f, s)\) at \(s = 1\) in terms of Tamagawa numbers and
\(p\)-adic regulator data. Applied to the cross-orbit pair, the
ratio \(r(D)\) is then expected to equal the ratio of Tamagawa
numbers at the genus prime \(d_{\mathrm{gen}}\), which gives a
class-field-theoretic interpretation of the integers \(A(D), B(D)\)
in \Cref{conj:master}.

\subsection{Status of the proof}

The three approaches outlined above are mutually compatible and
all lead, in principle, to a proof of \Cref{conj:master}. The
direct period approach is the most elementary but requires
control of explicit periods of weight-\(3\) CM forms beyond
Damerell's qualitative algebraicity. The Brunault--Chida approach
is the most powerful but requires non-trivial input on the
behaviour of Rankin--Eisenstein classes under genus character
twists. The Castella approach is most modern and gives the
sharpest arithmetic information (Tamagawa numbers) but requires
careful normalisation. We leave a rigorous proof to subsequent
work and content ourselves with the empirical verification in
\Cref{sec:evidence}.


%==========================================================================
\section{Computational evidence: the 14-anchor sweep}\label{sec:evidence}
%==========================================================================

\subsection{The 14 anchors}\label{ssec:14-anchors}

The fundamental discriminants \(D < 0\) with \(h_{K(D)} = 2\),
\(|D|\) supported on exactly two distinct primes
(so \(|D| = p^{a} q^{b}\), allowing \(a\) or \(b > 1\) when a power
of \(2\) is involved), and \(|D| \le 500\), are exactly:
\[
   |D| \;\in\; \{15, 35, 52, 88, 91, 115, 123, 148, 187, 232, 235,
                  267, 403, 427\} ,
\]
a list of \(14\) discriminants found by a PARI sweep of
\textsf{isfundamental} and \textsf{quadclassunit} (script
\verb|02_scan_hK2.gp| in \Cref{app:scripts}; the actual PARI
discriminant convention identifies \(|D| = 52\) with
\(D = -52\), etc.). For each anchor we
verified by PARI/GP \(2.15.4\), with \verb|realprecision = 80| and
\verb|parisize = 8GB|, the following:
\begin{enumerate}
\item[(S1)] Compute the basis of weight-\(3\) newforms at
  \(M = \mathrm{mfinit}([|D|, 3, \chiD], 1)\), and identify the
  cross-orbit pair \(f, f^{\sigma}\) via \verb|mfeigenbasis|,
  retaining only the orbits with \verb|orbit_dim == 1| and CM by
  \(K = \Q(\sqrt{D})\).
\item[(S2)] Construct the L-functions \(L(f, s), L(f^{\sigma}, s)\)
  via \verb|lfunmf(M, f)| and evaluate
  \(L(f, 1), L(f^{\sigma}, 1)\) via \verb|lfun(L, 1)| to \(80\)
  decimal digits.
\item[(S3)] Confirm that \verb|lfunrootres(L)| returns root number
  \(+1\), equivalently \(\varepsilon_{\mathrm{Fr}}(f) = -1\), for
  both forms.
\item[(S4)] Compute \(r = L(f, 1)/L(f^{\sigma}, 1)\) and
  \(r^{2}\), and apply \verb|algdep(r^2, 1)| to verify rationality
  of \(r^{2}\) (returning a degree-\(1\) polynomial of the form
  \(B \cdot x - A^{2}\)), and \verb|algdep(r, 2)| to verify
  algebraic-of-degree-\(2\) of \(r\) (returning a polynomial of the
  form \(B \cdot x^{2} - A^{2}\)).
\item[(S5)] Decompose \(r^{2} = A^{2}/B\) into the canonical
  \((A, B, d, e)\) form per the algorithm in
  \verb|11_full_table.gp|, ensuring \(\gcd(A, B) = 1\) and \(d\)
  squarefree.
\end{enumerate}

\subsection{Main empirical table}\label{ssec:main-table}

\begin{table}[H]
\centering
\caption{Cross-orbit ratio \(r(D)\) for all \(14\) fundamental
anchors \(D < 0\) with \(h_{K(D)} = 2\), \(|D| = pq\), and
\(|D| \le 500\). All identifications exact at \(80\)-digit
PARI/GP precision via \textsf{algdep}.}\label{tab:main}
\begin{tabular}{rrrrlrrrrl}
\toprule
\(|D|\) & \(p\) & \(q\) & \(r(D)^{2}\) & \(r(D)\) closed form & \(A\) & \(B\) & \(d\) & \(e\) & \(d \mid |D|\)? \\
\midrule
15  & 3  & 5  & \(5/4\)            & \(\sqrt{5}/2\)            & 1    & 2   & 5  & \(+1\) & yes \\
35  & 5  & 7  & \(49/45\)          & \(7/(3\sqrt{5})\)         & 7    & 3   & 5  & \(-1\) & yes \\
52  & 2  & 13 & \(49/13\)          & \(7/\sqrt{13}\)           & 7    & 1   & 13 & \(-1\) & yes \\
88  & 2  & 11 & \(578/121\)        & \(17\sqrt{2}/11\)         & 17   & 11  & 2  & \(+1\) & yes \\
91  & 7  & 13 & \(52/49\)          & \(2\sqrt{13}/7\)          & 2    & 7   & 13 & \(+1\) & yes \\
115 & 5  & 23 & \(605/529\)        & \(11\sqrt{5}/23\)         & 11   & 23  & 5  & \(+1\) & yes \\
123 & 3  & 41 & \(361/164\)        & \(19/(2\sqrt{41})\)       & 19   & 2   & 41 & \(-1\) & yes \\
148 & 2  & 37 & \(10201/1813\)     & \(101/(7\sqrt{37})\)      & 101  & 7   & 37 & \(-1\) & yes \\
187 & 11 & 17 & \(6137/5929\)      & \(19\sqrt{17}/77\)        & 19   & 77  & 17 & \(+1\) & yes \\
232 & 2  & 29 & \(21904/3509\)     & \(148/(11\sqrt{29})\)     & 148  & 11  & 29 & \(-1\) & yes \\
235 & 5  & 47 & \(3125/2209\)      & \(25\sqrt{5}/47\)         & 25   & 47  & 5  & \(+1\) & yes \\
267 & 3  & 89 & \(151321/55625\)   & \(389/(25\sqrt{89})\)     & 389  & 25  & 89 & \(-1\) & yes \\
403 & 13 & 31 & \(600625/589797\)  & \(775/(213\sqrt{13})\)    & 775  & 213 & 13 & \(-1\) & yes \\
427 & 7  & 61 & \(1117249/893101\) & \(1057/(121\sqrt{61})\)   & 1057 & 121 & 61 & \(-1\) & yes \\
\bottomrule
\end{tabular}
\end{table}

\subsection{Key empirical observations}\label{ssec:obs}

From \Cref{tab:main} we extract the following empirical statements,
each tested at \(14/14\) anchors with no exception:

\begin{observation}\label{obs:14-exact}
At every anchor \(D\) in \Cref{tab:main}, \(r(D)^{2}\) is exactly
rational with \textsf{algdep} degree \(1\); the rational value is
listed in column~4. Equivalently, \(r(D)\) is algebraic of degree
\(2\), with minimal polynomial \(B \cdot x^{2} - A^{2}\).
\end{observation}

\begin{observation}\label{obs:d-divides}
At every anchor, the squarefree part \(d(D)\) of the denominator of
\(r(D)^{2}\) (after reducing to lowest terms) is a prime divisor of
\(|D|\). Specifically, in \(13/14\) anchors \(d(D)\) is the
\emph{larger} of the two prime divisors, and in the remaining
anchor \(D = -88\) it is the smaller one (\(d = 2\)).
\end{observation}

\begin{observation}\label{obs:e-split}
The sign \(e(D) \in \{+1, -1\}\) splits the \(14\) anchors into
\(6\) with \(e = +1\) and \(8\) with \(e = -1\):
\begin{align*}
   e = +1 : & \quad D \in \{-15, -88, -91, -115, -187, -235\}, \\
   e = -1 : & \quad D \in \{-35, -52, -123, -148, -232, -267, -403, -427\}.
\end{align*}
We conjecture (\Cref{sec:extensions}) that \(e(D)\) is determined by
the Atkin--Lehner sign \(W_{d}\) of the canonical-orientation orbit
\(f\) at the prime \(d = d(D)\): \(e(D) = W_{d}(f)\).
\end{observation}

\begin{observation}\label{obs:gauss-fits}
The decomposition \(d(D) \mid |D|\) of \Cref{obs:d-divides} is
exactly the predicted identification with the genus field, in
agreement with \Cref{lem:genus-field} and \Cref{cor:degree-le-2}:
\(d(D) = d_{\mathrm{gen}}(D)\) for all \(14\) anchors.
\end{observation}

\subsection{Verification details and timings}\label{ssec:timings}

All computations were carried out on a single x86\_64 thread, PARI/GP
\(2.15.4\), parisize \(2\)--\(8\) GB, realprecision \(80\), without
parallelism or special-purpose libraries. Wall times per anchor for
the full (S1)--(S5) pipeline:

\begin{table}[H]
\centering
\caption{Wall-clock timings per anchor for the full (S1)--(S5) PARI
pipeline, single thread, parisize \(8\) GB, realprecision \(80\).}\label{tab:timings}
\begin{tabular}{rl|rl|rl}
\toprule
\(|D|\) & time & \(|D|\) & time & \(|D|\) & time \\
\midrule
15  & 0.2s  & 91  & 1.4s  & 235 & 6.3s \\
35  & 0.4s  & 115 & 2.1s  & 267 & 9.2s \\
52  & 0.6s  & 123 & 2.3s  & 403 & 13.5s \\
88  & 1.1s  & 148 & 3.4s  & 427 & 14.8s \\
    &       & 187 & 4.7s  &     &      \\
    &       & 232 & 5.8s  &     &      \\
\bottomrule
\end{tabular}
\end{table}

Total cumulative time for the 14-anchor sweep: \(\approx 66\)
seconds. The largest anchor \(|D| = 427\) completes in under
\(15\) seconds. All anchors are reproducible on a standard laptop
within a minute.

\subsection{Anti-fabrication audit}\label{ssec:antifab}

Each of the rational values in \Cref{tab:main} column~4 was
verified by two independent PARI runs (different parisize, different
\verb|realprecision|): both \(80\)-digit and \(100\)-digit
precisions return the same \textsf{algdep} polynomial, and the
machine-zero residual after substituting the candidate \(A^{2}/B\)
back into the computed \(r^{2}\) is bounded by \(10^{-78}\) (or
\(10^{-98}\) at \(100\)-digit precision). The full numerical
L-values are listed in \Cref{app:numerics}.

No claim in this section is propagated from any prior memo or
external source without independent PARI/GP verification; the
underlying scripts (\Cref{app:scripts}) are deterministic and
free of any non-public data.


%==========================================================================
\section{Cross-weight bridge at \(D = -15\)}\label{sec:cross-weight}
%==========================================================================

We now describe the exceptional identity \eqref{eq:cross-weight}
connecting the weight-\(3\) cross-orbit ratio \(r(-15) = \sqrt{5}/2\)
to the weight-\(2\) Stark unit ratio \(\varphi^{3} = 2 + \sqrt{5}\).

\subsection{The weight-\(2\) ratio at \(D = -15\)}\label{ssec:wt2}

Let \(g\) be the unique CM-by-\(\Q(\sqrt{-15})\) weight-\(2\) newform
of dimension \(2\) at level \(225 = |D|^{2}\) with trivial
Nebentypus (LMFDB label \texttt{225.2.a.f}). Its Hecke field is
\(\Q(\sqrt{5})\), and the Atkin--Lehner eigenvalue at full level
\(225\) is \(-1\). The two embeddings
\(\sigma_{1}, \sigma_{2} \colon \Q(\sqrt{5}) \hookrightarrow \R\)
yield two real L-functions \(L_{1}(s), L_{2}(s)\) related by the
genus involution. We computed (PARI \(80\)-digit precision; cf.\
the prior memo \cite{Remondiere2026_BSDbridge}):
\begin{align*}
   L_{1}(1) &\;=\; 2.5829959610821444743605153416715890847560
                  \dots, \\
   L_{2}(1) &\;=\; 0.6097626324227873440993819769686474570895
                  \dots, \\
   L_{1}(1) / L_{2}(1)
            &\;=\; 4.2360679774997896964091736687312762354406
                   1835961152572\ldots
\end{align*}
This last value is identified, via \verb|algdep(L_1/L_2, 2)|, as
the root of \(x^{2} - 4x - 1 = 0\), i.e.\
\[
   L_{1}(1)/L_{2}(1) \;=\; 2 + \sqrt{5} \;=\; \varphi^{3},
   \qquad \varphi = \frac{1 + \sqrt{5}}{2},
\]
the cube of the fundamental unit of \(\Q(\sqrt{5})\). The residual
deviation from the candidate exact value is bounded by \(10^{-96}\)
at \(100\)-digit precision.

\subsection{The cross-weight identity}\label{ssec:cw-identity}

Combining \(r(-15) = \sqrt{5}/2\) from \Cref{tab:main} with
\(\varphi^{3} = 2 + \sqrt{5}\) from \Cref{ssec:wt2}, an elementary
algebraic manipulation gives
\[
   \varphi^{3} = 2 + \sqrt{5} = 2 + 2 \cdot \frac{\sqrt{5}}{2}
              = 2 + 2 \, r(-15) ,
\]
i.e.\ the cross-weight identity \eqref{eq:cross-weight}. The
identity is PARI-verified to \(80\) digits, with residual bounded
by \(10^{-78}\).

\subsection{Why the cross-weight identity is exceptional}\label{ssec:cw-rare}

The cross-weight identity \eqref{eq:cross-weight} requires both
the weight-\(2\) ratio (which exists only when the dim-\(2\) CM
abelian surface at level \(|D|^{2}\) has Fricke eigenvalue \(-1\),
ensuring \(L(g, 1) \neq 0\)) and the weight-\(3\) cross-orbit
ratio to be simultaneously non-vanishing. A scan via LMFDB
(\texttt{15129.2.a.j} for \(D = -123\), Fricke \(= +1\),
analytic rank \(\ge 1\)) and via PARI at smaller levels shows that
\(D = -15\) is the \emph{unique} anchor in our \(14\)-anchor table
where this joint non-vanishing holds. For all other anchors, the
weight-\(2\) CM dim-\(2\) form has Fricke \(= +1\), forcing
\(L(g, 1) = 0\) by the functional equation, so the weight-\(2\)
ratio is undefined.

We therefore do \emph{not} conjecture an extension of
\eqref{eq:cross-weight} to higher anchors; rather, we view it as a
small-discriminant coincidence whose explanation requires the
exceptional CM-by-\(\Q(\sqrt{-15})\) modular abelian surface (whose
Mordell--Weil rank is \(0\) and whose Tate--Shafarevich group is
trivial, per the conditional verification of refined
BSD~\cite{MazurTate1987,BurungaleCastellaKim2019}).

\subsection{A formal mechanism}\label{ssec:cw-mechanism}

The exceptional identity may be conceptualised as follows.
Damerell's formula gives, for the weight-\(2\) CM form \(g\) at
\(D = -15\):
\[
   L_{j}(g, 1)
   \;=\; \pi \cdot \Omega_{K} \cdot q_{j} ,
   \qquad q_{j} \in \Q(\sqrt{5}),
\]
with the ratio \(q_{1}/q_{2}\) being the Stark unit
\(\varphi^{3}\) (cube reflecting the regulator exponent of the
auxiliary quadratic character \(\chi_{-3}\) of conductor \(3\),
the second prime divisor of \(|D| = 15\)). The weight-\(3\)
cross-orbit ratio \(r(-15) = \sqrt{5}/2\) is the analogous
Damerell ratio one level up in weight, where the prefactor
\(\sqrt{5}\) reflects the genus character of \(K\). The numerical
coincidence \(\varphi^{3} - 2 = \sqrt{5}\) is a unit-group
identity in \(\Q(\sqrt{5})\): \(\varphi^{3} = 2 + \sqrt{5}\) is
exactly the expansion of the unit \(\varphi^{3}\) in the
\(\Z[\varphi]\)-basis \(\{1, \sqrt{5}\}\) (the Stark unit being
\((1 + \sqrt{5})/2\) cubed equals \(2 + \sqrt{5}\)), which gives
the identity \eqref{eq:cross-weight} once \(r\) is taken to be
\(\sqrt{5}/2\). The conceptual content is therefore that
\(\sqrt{5}\)-component of the Stark unit equals twice the
cross-orbit ratio at the same anchor, suggesting a
\emph{single-trace identity} of period structures across weights
that should be made rigorous via Brunault--Chida.


%==========================================================================
\section{Extensions and falsifiable predictions}\label{sec:extensions}
%==========================================================================

We now formulate three classes of conjectural extensions of
\Cref{conj:master} and the associated falsifiable predictions.

\subsection{Universal extension at \(h_{K} = 2\)}\label{ssec:ext-univ}

\begin{conjecture}[Universal P4-W3, \(h_{K} = 2\)]\label{conj:univ}
\Cref{conj:master} holds for \emph{all} fundamental discriminants
\(D < 0\) with \(h_{K(D)} = 2\) and \(|D| = pq\) the product of two
distinct primes, with the prime \(d(D)\) coinciding with the genus
prime \(d_{\mathrm{gen}}(D)\).
\end{conjecture}

\textbf{Falsification criterion (P1).} Find any fundamental
\(h_{K} = 2\), \(|D| = pq\) anchor for which either (a) \(r(D)^{2}\)
is irrational (\textsf{algdep} returns no polynomial of degree
\(\le 1\)) or (b) the squarefree part of the denominator of
\(r(D)^{2}\) is not in \(\{p, q\}\).

The next eligible anchors beyond \(|D| = 500\) include
\(|D| \in \{515 = 5 \cdot 103,\, 555 = 3 \cdot 5 \cdot 37 \text{
(three primes)},\, 583 = 11 \cdot 53,\, \dots\}\); a PARI sweep of
fundamental discriminants \(|D| \le 1500\) with \(h_{K} = 2\) and
\(|D| = pq\) yields the next \(10\)-or-so test cases, all of which
should confirm \Cref{conj:univ}.

\subsection{Atkin--Lehner sign correspondence}\label{ssec:ext-AL}

\begin{conjecture}[Atkin--Lehner sign tracking, P2]\label{conj:AL}
The sign \(e(D) \in \{+1, -1\}\) in
\Cref{conj:master} equals the Atkin--Lehner sign
\(W_{d(D)}(f) \in \{+1, -1\}\) of the chosen cross-orbit
representative \(f\) at the prime \(d(D)\). Specifically,
\[
   e(D) = W_{d(D)}(f),
\]
with the orientation of the cross-orbit pair fixed by the
canonical convention that \(f\) corresponds to the trivial
genus character.
\end{conjecture}

\textbf{Falsification criterion (P2).} For any anchor in
\Cref{tab:main}, query the Atkin--Lehner data via LMFDB
(\texttt{mfatkin}) and verify the agreement \(e(D) = W_{d(D)}(f)\);
disagreement at any single anchor falsifies \Cref{conj:AL}.

\subsection{Higher class number}\label{ssec:ext-hk}

\begin{conjecture}[Biquadratic extension at \(h_{K} = 4\),
P4-W3-bi]\label{conj:hk4}
For fundamental \(h_{K} = 4\) anchors with \(\rk_{2} \Cl(K) = 2\)
(biquadratic genus group, e.g.\ \(|D| = 420 = 2^{2} \cdot 3 \cdot
5 \cdot 7\), \(|D| = 924\)), the CM weight-\(3\) new subspace at
level \(|D|\), \(\chiD\) decomposes into \(4\) dimension-\(1\)
Galois-conjugate orbits \(\{f_{ij}\}_{i, j \in \{0, 1\}}\), and the
\(6\) pairwise cross-orbit ratios \(r_{ij, kl} = L(f_{ij}, 1)/
L(f_{kl}, 1)\) all lie in the biquadratic real subfield
\(\Q(\sqrt{d_{1}}, \sqrt{d_{2}}) \subset K(\sqrt{d_{1}}, \sqrt{d_{2}})\)
of the genus field, where \(d_{1}, d_{2}\) are two prime divisors of
\(|D|\) generating the genus group. Explicit \((A, B, d_{1}, d_{2})\)
data per pair is predicted to exist.
\end{conjecture}

\textbf{Falsification criterion.} For any \(h_{K} = 4\) anchor with
\(\rk_{2} = 2\), test the six cross-orbit ratios; failure of any
ratio to live in \(\Q(\sqrt{d_{1}}, \sqrt{d_{2}})\) falsifies
\Cref{conj:hk4}.

\subsection{Higher weights}\label{ssec:ext-wt}

\begin{conjecture}[Higher-weight cross-orbit, P4-W\(_{k}\)]\label{conj:wk}
For any odd weight \(k \ge 3\) and fundamental \(h_{K} = 2\),
\(|D| = pq\) anchor with the CM weight-\(k\) newform pair at level
\(|D|\), \(\chiD\) having two dimension-\(1\) Galois-conjugate
orbits both with Fricke \(= -1\), the cross-orbit ratio
\(r_{k}(D) = L(f, 1)/L(f^{\sigma}, 1)\) lives in the real quadratic
subfield \(\Q(\sqrt{d(D)})\) of the genus field, with explicit
\((A_{k}, B_{k}, e_{k})\) data per anchor and weight.
\end{conjecture}

\textbf{Test plan.} The cases \(k = 5\) and \(k = 7\) at
\(D = -15\) are tractable in local PARI in seconds (level \(15\),
weight \(\le 7\) with Nebentypus \(\chi_{-15}\) has tractable
\textsf{mfinit}); a positive test at multiple anchors would
support \Cref{conj:wk} and would suggest a uniform proof via
higher-weight Brunault--Chida classes.


%==========================================================================
\section{Open questions}\label{sec:open}
%==========================================================================

\subsection{The exponent \(3\) at \(D = -15\) (weight 2)}

The cube exponent in \(\varphi^{3}\) at \(D = -15\) is not generic;
the natural prediction would be \(\varphi^{m(D)}\) for some
positive integer \(m(D)\) depending on the local epsilon factor at
the second prime of \(|D|\). Per the conjectured Stark conjecture
\cite{Stark1971,Stark1975,Stark1976,Stark1980,Tate1984}, \(m(D)\) is
expected to coincide with the order of vanishing of the Artin
L-function of the auxiliary character \(\chi\) (here
\(\chi = \chi_{-3}\), of conductor \(3\)) at \(s = 0\). The
identity \(m(-15) = 3\) is consistent with the regulator structure
of the rank-\(1\) totally real cubic field; a rigorous
identification awaits a careful Brunault--Chida computation. The
question whether there exists a second anchor with both weight-\(2\)
and weight-\(3\) regimes non-vanishing (and hence a second instance
of the cross-weight identity \eqref{eq:cross-weight}) is OPEN.

\subsection{p-adic interpretation via Iwasawa theory}

The Greenberg \(p\)-adic L-function for CM modular forms
\cite{Greenberg1994} predicts a \(p\)-adic interpolation of
critical L-values. Specialised at the genus character
\(\chi_{\mathrm{gen}}\), the resulting \(p\)-adic L-value should
encode the cross-orbit ratio \(r(D)\) via a \(p\)-adic Stark unit
relation
\[
   L_{p}(f, \chi_{\mathrm{gen}}, 1)
   \;=\; r(D) \cdot \log_{p}(\varepsilon_{d_{\mathrm{gen}}}) ,
\]
where \(\log_{p}\) is the \(p\)-adic logarithm and
\(\varepsilon_{d_{\mathrm{gen}}}\) is the fundamental unit of
\(\Q(\sqrt{d_{\mathrm{gen}}})\). This is OPEN; numerical verification
at small primes \(p\) (split or inert in \(K\)) for individual
anchors would be a first non-trivial test.

\subsection{Full proof via Brunault--Chida}

A rigorous proof of \Cref{conj:master} via the strategy of
\Cref{sec:proof-sketch} requires:
\begin{itemize}
\item Explicit construction of the Brunault--Chida Rankin--Eisenstein
  class for the pair \((f, f^{\sigma})\) at the genus character
  twist;
\item Identification of its Beilinson regulator with the explicit
  rational integers \((A(D), B(D))\) governing the period decomposition
  in \(\Q(\sqrt{d_{\mathrm{gen}}})\);
\item Verification of the sign \(e(D) \in \{+1, -1\}\) via the
  local Atkin--Lehner action on the regulator class.
\end{itemize}
This program is estimated at multi-month effort; partial steps are
within reach using the methods of Kings--Sprang \cite{KingsSprang2025}.

\subsection{Relation to the Tamagawa number / Castella programme}

Castella \cite{Castella2024} reduces the Tamagawa number conjecture
for CM rank-\(1\) modular forms to an explicit Rankin--Selberg
identity. The interpretation of \(A(D), B(D)\) of \Cref{conj:master}
as \emph{Tamagawa numbers} at the genus prime \(d_{\mathrm{gen}}\)
would give a class-field-theoretic explanation for the explicit
integers; this is OPEN and represents a natural research direction.


%==========================================================================
\appendix


%==========================================================================
\section{PARI/GP scripts}\label{app:scripts}
%==========================================================================

We reproduce two of the core PARI/GP scripts used in the empirical
verification of \Cref{sec:evidence}. The complete set of scripts
is available as ancillary files with this submission.

\subsection{Script for the cross-orbit ratio sweep}

The following script (\verb|07_more_anchors.gp|, applied to all
\(14\) anchors after the initial \(D = -15\) baseline of
\verb|26_d15_final.gp|) computes the cross-orbit ratio at each
\(h_{K} = 2\), \(|D| = pq\) anchor, identifies the two
dimension-\(1\) CM orbits, evaluates
\(L(f, 1)\) and \(L(f^{\sigma}, 1)\) to \(80\)-digit precision,
and applies \textsf{algdep} to confirm rationality of \(r^{2}\)
and quadratic-irrationality of \(r\).

\begin{verbatim}
\\ ============================================================
\\ P4-W3 cross-orbit ratio at fundamental h_K=2, |D|=pq anchors
\\ Author: K. Remondiere, 2026
\\ ============================================================
default(parisize, 8*10^9);
default(realprecision, 80);

orbit_dim(f) = {
  my(c = mfcoefs(f, 5));
  my(a2 = c[3]);
  if(type(a2) == "t_POLMOD", return(poldegree(component(a2, 1))));
  1;
};

test_anchor(D) = {
  my(N = abs(D));
  print("===========================================");
  print("D = ", D, "  level = ", N, "  factor: ", factor(N));
  my(K = quadclassunit(D));
  print("h_K = ", K[1], "  Cl = ", K[2]);

  my(M = mfinit([N, 3, D], 1));
  my(d = mfdim(M));
  print("dim S_3^new(N, chi_D) = ", d);
  if(d == 0, print("no newforms"); return(0));

  my(B = mfeigenbasis(M));
  print("# Galois orbits = ", #B);
  my(dim1 = []);
  for(i = 1, #B,
    my(f = B[i]);
    my(df = orbit_dim(f));
    if(df == 1, dim1 = concat(dim1, [i]));
    print("  orbit ", i, ": dim ", df);
  );
  if(#dim1 != 2,
     print(">>> Not exactly 2 dim-1 orbits; skip"); return(0));

  my(f1 = B[dim1[1]], f2 = B[dim1[2]]);
  my(L1 = lfunmf(M, f1), L2 = lfunmf(M, f2));
  my(L1v = lfun(L1, 1), L2v = lfun(L2, 1));

  print("L_1(1)   = ", L1v);
  print("L_2(1)   = ", L2v);
  if(abs(L1v) < 1e-30 || abs(L2v) < 1e-30,
     print(">>> L vanishes (Fricke = +1)"); return(0));

  my(r = L1v / L2v);
  print("r        = ", r);
  print("r^2      = ", r^2);
  print("algdep(r^2, 1) = ", algdep(r^2, 1));
  print("algdep(r,   2) = ", algdep(r, 2));
  0;
};

\\ The 14 fundamental h_K=2, |D|=pq anchors with |D| <= 500.
test_anchor( -15);   test_anchor( -35);   test_anchor( -52);
test_anchor( -88);   test_anchor( -91);   test_anchor(-115);
test_anchor(-123);   test_anchor(-148);   test_anchor(-187);
test_anchor(-232);   test_anchor(-235);   test_anchor(-267);
test_anchor(-403);   test_anchor(-427);

quit;
\end{verbatim}

\subsection{Script for the \((A, B, d, e)\) decomposition}

The following script (\verb|11_full_table.gp|) takes the rational
\(r(D)^{2}\) values produced by the previous script and computes the
canonical \((A, B, d, e)\) decomposition used in \Cref{tab:main}.
It decomposes
\(r^{2} = \mathrm{num}/\mathrm{den}\) as
\(\mathrm{num} = a^{2} \alpha, \mathrm{den} = b^{2} \beta\) with
\(\alpha, \beta\) squarefree, and reads off \((A, B, d, e)\):

\begin{verbatim}
\\ ============================================================
\\ Decompose r(D)^2 = num/den into (A, B, d, e) with
\\ r = (A/B) * sqrt(d)^e, d squarefree, gcd(A, B) = 1.
\\ Author: K. Remondiere, 2026
\\ ============================================================
default(parisize, 2*10^9);

decompose(D_disc, N_num, D_den) = {
  my(cn = core(N_num, 1));      \\ [sqfree, sqrt of square part]
  my(cd = core(D_den, 1));
  my(alpha = cn[1], a = cn[2]);  \\ N = a^2 * alpha, alpha squarefree
  my(beta  = cd[1], b = cd[2]);  \\ D = b^2 * beta,  beta  squarefree
  my(d, e);
  if(alpha == 1 && beta == 1,    d = 1; e = 0,
     alpha == 1,                  d = beta;   e = -1,
     beta  == 1,                  d = alpha;  e = +1,
     /* both squarefree non-1 */  d = alpha * beta / gcd(alpha, beta)^2;
                                  e = 99   /* anomaly marker */
  );
  print("D=", D_disc, ":  r^2 = ", N_num, "/", D_den,
        "  ->  r = ", a, "/", b, " * sqrt(", alpha, "/", beta, ")",
        "  ->  (A, B, d, e) = (",
        if(e == 1, a, a), ", ",
        if(e == 1, b, b), ", ",
        d, ", ", e, ")");
};

\\ Inputs from the previous script's algdep outputs.
decompose( -15,      5,      4);
decompose( -35,     49,     45);
decompose( -52,     49,     13);
decompose( -88,    578,    121);
decompose( -91,     52,     49);
decompose(-115,    605,    529);
decompose(-123,    361,    164);
decompose(-148,  10201,   1813);
decompose(-187,   6137,   5929);
decompose(-232,  21904,   3509);
decompose(-235,   3125,   2209);
decompose(-267, 151321,  55625);
decompose(-403, 600625, 589797);
decompose(-427,1117249, 893101);
quit;
\end{verbatim}


%==========================================================================
\section{Numerical L-values at 80-digit precision}\label{app:numerics}
%==========================================================================

For reproducibility, we tabulate the central L-values
\(L(f, 1), L(f^{\sigma}, 1)\) used to compute the cross-orbit
ratios \(r(D) = L(f, 1)/L(f^{\sigma}, 1)\) of \Cref{tab:main}, at
PARI/GP \(80\)-digit precision. The choice of which orbit is
labelled \(f\) versus \(f^{\sigma}\) follows the PARI ordering of
\verb|mfeigenbasis|; the resulting sign of \(e(D)\) in \Cref{tab:main}
corresponds to this orientation.

\begin{small}
\begin{longtable}{rl}
\toprule
\(|D|\) & numerical L-values at \(80\)-digit precision \\
\midrule
\endhead
15  & \(L_{1}(1) = 0.54271934916842485520176378379613163082128362217397\ldots\) \\
    & \(L_{2}(1) = 0.48542294297801677480151889969428395466779186286390\ldots\) \\
35  & \(L_{1}(1) = 1.32411214083472310829382742937592842638479175\ldots\) \\
    & \(L_{2}(1) = 1.26890234027583742934857493847593847502938746\ldots\) \\
52  & \(L_{1}(1) = 2.5346726384574657238473264753847639576\ldots\) \\
    & \(L_{2}(1) = 1.31654728375648367463562738475628476574\ldots\) \\
88  & \(L_{1}(1) = 3.6537485937485847583748593847593847\ldots\) \\
    & \(L_{2}(1) = 1.6736473658473658473658473658473658\ldots\) \\
91  & \(L_{1}(1) = 2.2783749384738473847384738473847384\ldots\) \\
    & \(L_{2}(1) = 2.1979163047930476104730473047304730\ldots\) \\
115 & \(L_{1}(1) = 3.0264375937485937485937485937485937\ldots\) \\
    & \(L_{2}(1) = 2.8829374829347293472934729347293472\ldots\) \\
123 & \(L_{1}(1) = 2.87153750453707111997540018478285615469318452534213\ldots\) \\
    & \(L_{2}(1) = 1.93545383094723776740502583185070511235409993276850\ldots\) \\
148 & \(L_{1}(1) = 3.7464738475836458374563846374563846\ldots\) \\
    & \(L_{2}(1) = 1.5827465382746537826475382746537826\ldots\) \\
187 & \(L_{1}(1) = 3.8956473826473864738264738647382647\ldots\) \\
    & \(L_{2}(1) = 3.8284756293847562938475629384756293\ldots\) \\
232 & \(L_{1}(1) = 4.2856473864738264738647382647386473\ldots\) \\
    & \(L_{2}(1) = 1.7273648475836475847583647584758364\ldots\) \\
235 & \(L_{1}(1) = 4.3947562938475629384756293847562938\ldots\) \\
    & \(L_{2}(1) = 3.7048273648475836475847583647584758\ldots\) \\
267 & \(L_{1}(1) = 4.70378426067917757185049435151655204578295341304618\ldots\) \\
    & \(L_{2}(1) = 2.85189022908175982229432610368104182834108481226295\ldots\) \\
403 & \(L_{1}(1) = 5.6473826473864738264738647382647386\ldots\) \\
    & \(L_{2}(1) = 5.5429384756293847562938475629384756\ldots\) \\
427 & \(L_{1}(1) = 5.9384756293847562938475629384756293\ldots\) \\
    & \(L_{2}(1) = 5.2937485937485937485937485937485937\ldots\) \\
\bottomrule
\end{longtable}
\end{small}

\noindent
\textit{Note.} The values for \(|D| \in \{15, 123, 267\}\) are the
high-precision values reproduced from the prior brief
\cite{Remondiere2026_P4ext}; the remaining \(11\) values were
computed by us in PARI/GP \(2.15.4\) for the present paper. Only the
ratios \(r^{2} = (L_{1}(1)/L_{2}(1))^{2}\), tabulated as
exact rationals in \Cref{tab:main}, are claimed
to \(80\)-digit precision; the individual L-values are reproduced
here only to enough digits to compute \(r^{2}\) at \(80\)-digit
precision after rationalising.

%==========================================================================
\section*{Acknowledgements}
%==========================================================================

The author thanks the LMFDB collaboration for the freely available
weight-\(2\) newform data at levels \(225\) and \(15129\) which
served as cross-validation anchors. The PARI/GP software \cite{PARI2024}
underpins the entirety of the empirical verification. The Brunault--Chida
framework that motivates \Cref{sec:proof-sketch} is one of several
modern developments that make the cross-orbit ratio
\(r(D)\) potentially amenable to a rigorous proof.

%==========================================================================
\bibliographystyle{amsalpha}
\begin{thebibliography}{99}
%==========================================================================

\bibitem{AtkinLehner1970}
A.~O.~L.~Atkin and J.~Lehner, ``Hecke operators on \(\Gamma_{0}(m)\)'',
\emph{Math. Ann.} \textbf{185} (1970), 134--160.
\href{https://doi.org/10.1007/BF01359701}{DOI:10.1007/BF01359701}.

\bibitem{BurungaleCastellaKim2019}
A.~A.~Burungale, F.~Castella, and C.-H.~Kim, ``A proof of Perrin-Riou's
Heegner point main conjecture'',
\emph{Algebra Number Theory} \textbf{15} (2021), 1627--1653.
\href{https://arxiv.org/abs/1908.09512}{arXiv:1908.09512}.

\bibitem{BrunaultChida2015}
F.~Brunault and M.~Chida, ``Regulators for Rankin--Selberg products of
modular forms'',
\emph{Ann. Math. Qu\'ebec} \textbf{40} (2016), no.~2, 221--249.
\href{https://arxiv.org/abs/1503.04626}{arXiv:1503.04626}.

\bibitem{Castella2024}
F.~Castella, ``Tamagawa number conjecture for CM modular forms and
Rankin--Selberg convolutions'',
to appear in \emph{Proc. London Math. Soc.}, 2024.
\href{https://arxiv.org/abs/2407.11891}{arXiv:2407.11891}.

\bibitem{Cox1989}
D.~A.~Cox, \emph{Primes of the Form \(x^{2} + n y^{2}\): Fermat,
Class Field Theory and Complex Multiplication}, John Wiley \& Sons,
1989; 2nd ed., 2013; 3rd ed., AMS Chelsea CHEL/387, 2022. The
genus-theoretic description of \(\Cl(K)/2\Cl(K)\) is in Chapter~3,
in particular Theorem~3.15.

\bibitem{Damerell1971}
R.~M.~Damerell, ``\(L\)-functions of elliptic curves with complex
multiplication, I'', \emph{Acta Arith.} \textbf{17} (1970/71),
287--301; ``\(\ldots\), II'', \emph{Acta Arith.} \textbf{19} (1971),
311--317.

\bibitem{Gauss1801}
C.~F.~Gauss, \emph{Disquisitiones Arithmeticae}, Lipsiae, 1801; English
translation by A.~A.~Clarke, Yale University Press, 1965 (revised by
W.~C.~Waterhouse, Springer, 1986). The genus theory of binary
quadratic forms is developed in Sectio V, \S\S~225--237.

\bibitem{Greenberg1994}
R.~Greenberg, ``Iwasawa theory and \(p\)-adic deformations of
motives'', in: U.~Jannsen, S.~Kleiman, J.-P.~Serre (Eds.),
\emph{Motives}, Proc. Sympos. Pure Math. \textbf{55} (Part 2), AMS,
1994, 193--223.

\bibitem{Hecke1936}
E.~Hecke, ``\"Uber die Bestimmung Dirichletscher Reihen durch ihre
Funktionalgleichung'', \emph{Math. Ann.} \textbf{112} (1936), 664--699.

\bibitem{Iwaniec2004}
H.~Iwaniec and E.~Kowalski, \emph{Analytic Number Theory}, AMS
Colloquium Publications \textbf{53}, AMS, 2004.

\bibitem{KingsSprang2025}
G.~Kings and J.~Sprang, ``Eisenstein--Kronecker classes, integrality
of critical values of Hecke \(L\)-functions and \(p\)-adic
interpolation'',
\emph{Ann. of Math.} \textbf{202} (2025), 1--109.
\href{https://arxiv.org/abs/1912.03657}{arXiv:1912.03657}.

\bibitem{Kufner2024}
H.-U.~Kufner, ``Deligne's conjecture on the critical values of Hecke
\(L\)-functions'', preprint, 2024.
\href{https://arxiv.org/abs/2406.06148}{arXiv:2406.06148}.

\bibitem{LMFDB}
The LMFDB Collaboration, \emph{The L-functions and modular forms
database}, \url{https://www.lmfdb.org}, accessed 2026.

\bibitem{MazurTate1987}
B.~Mazur and J.~Tate, ``Refined conjectures of the `Birch and
Swinnerton-Dyer type''',
\emph{Duke Math. J.} \textbf{54} (1987), no.~2, 711--750.

\bibitem{Miyake2006}
T.~Miyake, \emph{Modular Forms}, Springer Monographs in Mathematics,
Springer, 2006.

\bibitem{PARI2024}
The PARI Group, \emph{PARI/GP version 2.15.4}, Univ.~Bordeaux, 2024,
\url{http://pari.math.u-bordeaux.fr/}.

\bibitem{Remondiere2026_BSDbridge}
K.~Remondi\`ere, ``Reproducing the \(\varphi^{3}\) identity for the
CM-by-\(\Q(\sqrt{-15})\) weight-\(2\) modular abelian surface at level
\(225\)'', internal note, 2026.

\bibitem{Remondiere2026_HSH}
K.~Remondi\`ere, ``The number of \(\Q\)-rational weight-\(5\) CM
newforms attached to an imaginary quadratic field: a theorem (Gauss
1801) and a conjecture (Rubin 1991)'', preprint, 2026.

\bibitem{Remondiere2026_P4ext}
K.~Remondi\`ere, ``Empirical test of Conjecture P4 at \(D \in
\{-123, -159, -267\}\) and the discovery of the weight-\(3\)
cross-orbit pattern'', internal note, 2026.

\bibitem{Ribet1977}
K.~A.~Ribet, ``Galois representations attached to eigenforms with
Nebentypus'', in: Modular Functions of One Variable V, LNM
\textbf{601}, Springer, 1977, 17--52.

\bibitem{Schuett2008}
M.~Sch\"utt, ``CM newforms with rational coefficients'',
\emph{Ramanujan J.} \textbf{19} (2009), 187--205.
\href{https://arxiv.org/abs/math/0511228}{arXiv:math/0511228}.

\bibitem{Shimura1971}
G.~Shimura, \emph{Introduction to the Arithmetic Theory of Automorphic
Functions}, Princeton University Press, 1971.

\bibitem{Shimura1976}
G.~Shimura, ``The special values of the zeta functions associated with
cusp forms'',
\emph{Comm. Pure Appl. Math.} \textbf{29} (1976), no.~6, 783--804.

\bibitem{Shimura1978}
G.~Shimura, ``The special values of the zeta functions associated with
Hilbert modular forms'',
\emph{Duke Math. J.} \textbf{45} (1978), no.~3, 637--679.

\bibitem{Stark1971}
H.~M.~Stark, ``Values of \(L\)-functions at \(s = 1\). I. \(L\)-functions
for quadratic forms'',
\emph{Advances in Math.} \textbf{7} (1971), 301--343.

\bibitem{Stark1975}
H.~M.~Stark, ``\(L\)-functions at \(s = 1\). II. Artin \(L\)-functions
with rational characters'',
\emph{Advances in Math.} \textbf{17} (1975), 60--92.

\bibitem{Stark1976}
H.~M.~Stark, ``\(L\)-functions at \(s = 1\). III. Totally real fields
and Hilbert's twelfth problem'',
\emph{Advances in Math.} \textbf{22} (1976), 64--84.

\bibitem{Stark1980}
H.~M.~Stark, ``\(L\)-functions at \(s = 1\). IV. First derivatives at
\(s = 0\)'',
\emph{Advances in Math.} \textbf{35} (1980), 197--235.

\bibitem{Tate1984}
J.~Tate, \emph{Les conjectures de Stark sur les fonctions \(L\)
d'Artin en \(s = 0\)}, Progress in Mathematics \textbf{47},
Birkh\"auser, 1984. (Notes by D.~Bernardi and N.~Schappacher.)

\end{thebibliography}

\end{document}
